% New method for one ideal fourth integral; no frozen-tree update.
\subsection{Affine parent bounds and a concave-corner fourth integral}

Fix the same causal ideal policy $C_3=\Phi_3$ and the maximum $C_4$
over all licensed fourth selected-subproduct caps and its inherited
parent, at $w=19/125$. We bound
\[
 I_4=4\int_{E_4}
 \frac{\bar f((C_4-S_3-r)/r)-\bar f((C_3-S_3-r)/r)}
 {t_1t_2t_3r^2}\,d\boldsymbol t\,dr.
\]
All omitted contributions are nonnegative. Bounds in this note are
alternative lower estimates of $I_4$, not additions to previous
fourth totals.

\paragraph{A symmetric representation of the parent policy.}
On the ordered retained domain, $t_j\le w<25/128$ and
$t_1+3t_2<\nu$. The reduced eight-branch equivalence gives
\[
 C_3=\max\left\{\nu,
    \max_{\varnothing\ne I\subset\{1,2,3\}}v\Bigl(\sum_{i\in I}t_i\Bigr),
    \max_{1\le j\le3}\left(\Lambda-\frac38t_j\right)\right\}.
\]
The additional singleton $v(t_3)$ is dominated by $v(t_1)$;
the extra $W$ terms are dominated by the one for the smallest
coordinate. Thus this exactly represents the ordered cap and
extends it symmetrically to unordered boxes by sorting.
For every nonnegative argument the middle affine formula in $v$
is redundant, so
$v(s)=\min\{(1+s)/2,57/64-3s/2\}$.
Indeed the two outer lines cross at $s=25/128$, where the middle
line is strictly higher; comparing their slopes on the two sides
proves the assertion on the entire nonnegative half-line.

\paragraph{A genuine affine majorant of the old cap.}
Write each parent branch as $B_k(t)=\min_{q\in\mathcal L_k}q(t)$,
where singleton lists include the constant and $W$ branches.
Choose an affine plane $P$ from a branch active at an approximate
logarithmic mean. Put
\[
 \delta=\max\left\{0,
     \max_k\min_{q\in\mathcal L_k}\max_{t\in B}(q(t)-P(t))\right\},
 \qquad U(t)=P(t)+\delta.
\]
For every $k$ and every $t\in B$, the minimum defining $B_k(t)$
is at most each $q(t)$, so $B_k(t)\le U(t)$. Hence $C_3\le U$
throughout the whole box. Every displayed extremum is affine and
is evaluated at its exact signed endpoint. A constant parent upper
bound may be used instead if it gives a smaller value at the chosen
mean. The approximate mean affects this choice only, not validity.
A zero shift alone is not claimed to imply equality $U=C_3$ across
a parent branch wall.

The historical cube always uses a separate lower bound $L_3\le C_3$:
\[
 t_{1,+}+3t_{2,+}<\nu,\qquad
 0<r<\min\{t_{3,-},(L_3-\textstyle\sum t_{j,+})/3\}.
\]
The majorant $U$ is solely an integration device for the subtracted
old density. It is not substituted for the actual parent's cap in
the historical gate or the causal tree.

\paragraph{Whole-box licensed new endpoints.}
The estimator considers masks $9,10,14$, all including the last
variable $r$. For a fixed physical subset of the old coordinates,
the candidate is $V(t,r)=v(r+\sum_{j\in I}t_j)$.
Both affine pieces are tested against every selected-remaining
square gate throughout the entire outer box and retained $r$ interval.
Where outer bins coincide, every possible order of the unselected
old coordinates is tested. Thus each candidate is a valid lower
bound for the same full ideal $C_4$, even when sorting changes its
numerical mask label.

Also require, using exact affine extrema,
$V(t,r)-\sum t_j>3r$ throughout. This keeps the new ordinary argument
strictly above two, the range needed for concavity. It is harmless
to omit points failing this additional numerical-estimator restriction.
There is no requirement that $V>U$: the resulting lower gap is
allowed to be signed. The actual ideal fourth gap remains nonnegative.

\paragraph{Jensen for the old term and corners for the new term.}
On the outer box $B=\prod[a_j,b_j]$, use the independent logarithmic
probability measures $dt_j/(t_j\log(b_j/a_j))$ and let $m_j=Et_j$.
Since $\bar f$ is increasing and concave on $[2,\infty)$,
\[
 t\longmapsto -\bar f((U(t)-\sum t_j-r)/r)
\]
is convex. Jensen therefore bounds its expectation below by its
value at the true logarithmic mean.

The function $V(t,r)-\sum t_j$ is concave in the three old variables.
Composition with the increasing concave $\bar f$ gives a concave
new-density function on the whole retained box. Consequently its
multilinear corner interpolant is a pointwise lower bound. Taking
expectations yields
\[
 E\bar f((V-S_3-r)/r)
 \ge\sum_{\epsilon\in\{0,1\}^3}
 p_\epsilon\,
 \bar f((V(t^\epsilon,r)-\sum t_j^\epsilon-r)/r),
\]
where
\[
 p_\epsilon=\prod_{j=1}^3
 \begin{cases}
 (b_j-m_j)/(b_j-a_j),&\epsilon_j=0,\\
 (m_j-a_j)/(b_j-a_j),&\epsilon_j=1.
 \end{cases}
\]
All eight weights are nonnegative. This lower bound crosses the
tent-shaped $v$ kink without any derivative or curvature assumption.

Mean intervals and corner weights are enclosed outward. In the old
term use an upper endpoint for
$E(U-S_3)=U(m)-\sum m_j$; monotonicity of $\bar f$ gives the correct
direction for the subtracted integral. In the new term use lower
weights and lower nonnegative integrals. Selection among candidate
masks can use approximate mean values, since any selected licensed
candidate provides a valid lower bound. Interval partitions and
chosen candidates are fixed before integrating.

\paragraph{Exact last-variable integration and fallback.}
Put $H(u)=u\bar F(u)$, so $H'(u)=\bar f(u-1)$ for $u\ge3$.
At each corner the candidate $v$ has only its two affine pieces;
split at their exact crossing. For an affine last cap $A+Br$,
\[
 \int_\alpha^\beta \bar f((A+Br-S-r)/r)\frac{dr}{r^2}
 =\frac{H((A-S)/\alpha+B)-H((A-S)/\beta+B)}{A-S}.
\]
The old mean cap is handled with $B=0$. The positive logarithmic
outer mass is $4\prod_j\log(b_j/a_j)$. Table-range cutoffs are
exact, and small omitted $r$ tails are not assigned any gain.

Subtract the old upper integral from the weighted corner lower
integrals on each disjoint last-variable piece, and clip the result
at zero using nonnegativity of the true ideal gap. Independently
compute the previously licensed constant-corner estimate on the
same outer parent. Take the maximum of the two complete parent
lower bounds, never their sum. This preserves the earlier estimate
when the affine or corner enclosure is less effective.

\paragraph{Outer ordering, old history and the last ordering strip.}
If an outer interval occurs with multiplicity $m$, sorting divides
its full product box into $m!$ equal-measure sectors for the symmetric
actual ideal integrand. Divide the full-box lower integral by the
product of these factorials. The selected numerical bound need not
itself be symmetric. Gate checks for every possible remaining-slot
ordering ensure it is still a lower bound on every sector.

Old-history boundary cells are bisected only to a fixed finite depth.
Bisecting a group of $m$ equal intervals gives the $m+1$ ordered
choices of high/low multiplicities. Their factorial weights give
exactly the disjoint sector partition. Only cells whose whole old
history is valid are retained; unresolved leaves are omitted.

When the lowest outer interval $[a,b]$ is unique, restore part of
$a<r<t_3<b$. Write
$u=\min\{b,(L_3-\sum t_{j,+})/3\}$ and require $u>a$.
Its exact $(t_3,r)$ kernel mass is
\[
 \int_a^u\frac{dr}{r^2}\int_r^b\frac{dt_3}{t_3}
 =\frac{\log(b/a)}a-\frac{\log(b/u)}u-\frac1a+\frac1u.
\]
A whole-box constant new lower cap exceeding the old upper cap,
all selected gates and the ordinary argument constraints gives a
common-corner gap at $(t_{1,-},t_{2,-},a,a)$. Multiply it by this
mass and the other two logarithmic widths. This strip is disjoint
from the analytic slices $r<a$. Repeated-lowest-bin last strips and
the remaining variable-history boundary region are omitted.

\paragraph{Scope.}
The resulting total is a strict lower estimate for the same ideal
fourth integral, using finite rational partitions and the existing
outward ordinary table. It changes neither $C_3$ nor the all-mask
$C_4$ policy used in the fifth integral. A full arithmetic consequence
still requires one common finite causal approximation and the
previous compiler/distribution inputs. Numerical output and a
separate independent replay must be reported with their actual scope.
